\documentclass[12pt]{hw}
\usepackage{graphicx}
\usepackage{float}


%====================================================================
% IMPORTANT INFORMATION TO FILL OUT:
%====================================================================
\usepackage{amsthm}
\newtheorem{definition}{Definition}
\newcommand{\getr}{\ensuremath{{\overset{\$}{\leftarrow}}}\xspace}
\newcommand{\Gen}{\ensuremath{{\sf Gen}}\xspace}
\newcommand{\D}{\mathcal{D}}
\newcommand{\Z}{\mathbb{Z}}
\newcommand{\G}{\mathbb{G}}
\newcommand{\Dec}{\ensuremath{{\sf Dec}}\xspace}
\newcommand{\Enc}{\ensuremath{{\sf Enc}}\xspace}
% ENTER YOUR NAME HERE:
\renewcommand{\student}{ YOUR NAME HERE (uniqname here) }

% LIST ALL COLLABORATORS HERE:
\renewcommand{\collaborators}{Name/Exercise Number(s), Name/Exercise Numbers (s)}

%====================================================================
%====================================================================

\setcounter{sheet}{5} % Homework number here

\begin{document}

\maketitle

\begin{center} Due: Nov 26, 2024 (at 11:59 PM US Eastern Time) \\[10pt]
  \hrule height 1pt
\end{center}

\toggletrue{answers}
% \togglefalse{answers}

\setcounter{MaxMatrixCols}{20} \setlength\parindent{0pt}

% =============================================================================
% =============================================================================


% ===============================================
% ===============================================

\begin{exercise}
{\bf (15 points)}
In class, we saw how 
Ke can prove that she knows the password to a bike chain lock in zero-knowledge.
Now, suppose that Ke brings two bike chain locks. She wants to prove
in zero-knowledge 
that she knows the password for at least one of them, but without
revealing the password, or revealing which one she knows the password for.
Can you design a physical zero-knowledge proof protocol
such that Ke can prove the statement in zero knowledge?
Please describe the protocol clearly; you can draw a simple picture if that helps. 
You don't have to prove the security of the protocol.
\end{exercise}
\begin{answer}
Write Solution Here.
\end{answer}
% =============================================================================


\begin{exercise}
     {\bf (20 points)} 
     Suppose you are given two commitment schemes $(\mathsf{com}_1, \mathsf{open}_1)$ and $(\mathsf{com}_2, \mathsf{open}_2)$ with the same message space and commitment space.


\begin{enumerate}[(a)]
    \item Suppose that both commitment schemes satisfy the perfectly binding property. However only one of them satisfies the computationally hiding property (but you don't know which one). 
    Construct a commitment scheme $(\mathsf{com}, \mathsf{open})$ which is both perfectly binding and computationally hiding using only the provided commitment schemes.

    \item Now suppose that both the commitments satisfy the computationally hiding property. However, only one of them satisfies the perfectly binding property (but you don't know which one). Construct a commitment scheme $(\mathsf{com}, \mathsf{open})$ which is both perfectly binding and computationally hiding using only the provided commitment schemes. 
\end{enumerate}
\end{exercise}

\begin{answer}
  Write Solution Here.
\end{answer}


% =============================================================================
\begin{exercise}
{\bf (25 points)}
{\bf Graph Isomorphism}
Consider two graphs $G_0 = (V_0,E_0) $ and $G_1= (V_1,E_1)$ where $V_b$ and $E_b$ are the vertices and edges respectively for graph $G_b$. We say the graphs are isomorphic if they are structurally the same. Formally, $G_0$ and $G_1$ are isomorphic if there exists a bijective function $\varphi: V_0 \to V_1$ such that for any two vertices $u,v \in V_0$, $(u,v) \in E_0$ if and only if $(\varphi(u),\varphi(v)) \in E_1$. Such a function $\varphi$ is called an isomorphism.
Paxton wants to prove to Valerie that $G_0$ is isomorphic to $G_1$ (denoted $G_0 \simeq G_1$),
but does not want to reveal any information about $\varphi$.
Consider the following protocol.
\vspace{1em}
\begin{mdframed}
    \vspace{0.5em}
    \textbf{Zero-Knowledge Proof for Graph Isomorphism}\\\\
    Initially, Paxton takes as input statement $(G_0, G_1)$ and witness $\varphi$. Valerie takes as input statement $(G_0, G_1)$.
    Paxton and Valerie repeat the following interaction for $n$ rounds where $n$ is the security parameter.
    \begin{enumerate}
        \item Paxton picks a random permutation $\pi$, computes graph $H = \pi(G_1)$, and sends $H$ to Valerie. Here, $\pi(G_1)$ essentially relabels the vertices of the graph based on the permutation.
        \item Valerie responds with a challenge $c\ \sample\ \{0, 1\}$.
        \item If $c = 0$, Paxton sends $\rho = \pi \circ \varphi$, and Valerie checks if $H = \rho(G_0)$. Here $\circ$ denotes functional composition (i.e. $\rho(G)$ first computes $G' = \varphi(G)$ then computes $\pi(G')$)
        \item If $c = 1$, Paxton sends $\rho = \pi$, and Valerie checks if $H = \rho(G_1)$.
    \end{enumerate}
\end{mdframed}
\vspace{1em}


\begin{enumerate}[(a)]
\item Suppose that $G_0$ is not isomorphic to $G_1$. 
Show that for any possibly malicious prover, the probability 
that 
Valerie rejects in a \emph{single} fixed round is at least $\frac12$. 
Please justify your answer.

\item 
Suppose that $G_0$ is not isomorphic to $G_1$, and the prover 
can be arbitrarily malicious.
Please provide a lower bound on the probability that Valerie rejects in at least
one out of $n$ rounds. Give the tightest 
lower bound you can show, and justify your answer.

\item 
Prove that the above construction is honest-verifier zero-knowledge. 
Intuitively, 
you want to show that the transcript the verifier sees by interacting with an honest prover is indistinguishable from a simulated transcript generated without knowledge of the witness.
We can capture this intuition by showing there exists a simulator that can spoof an accepting proof for a valid statement, even without the corresponding witness, so long as it sees the verifier's randomness ahead of time.

Formally, if indeed $G_0 \simeq G_1$, there must exist a PPT simulator $\mathcal{S}$ that on input $(G_0, G_1)$ can produce a message transcript $(H', c', \rho')$ indistinguishable from one generated by Paxton and Valerie. You may assume that Valerie always behaves honestly.
Please construct such a simulator. 
Show that your simulator is PPT, and informally argue why it produces a transcript indistinguishable from one generated by Paxton and Valerie.


\end{enumerate}


\end{exercise}
\begin{answer}
Write Solution Here.
\end{answer}
% =============================================================================


\begin{exercise}
{\bf (25 points)}
{\bf Multi-Valued Byzantine Broadcast}
Consider a multi-valued variant of the original Byzantine Broadcast protocol (called Multi-Valued Byzantine Broadcast) in which the sender has an $\ell$-bit value ${\bf m} \in \{0, 1\}^\ell$, and it wants to propagate this value to everyone else. Consistency requires that all honest output the same value; and validity requires that if the sender is honest, all honest nodes output the sender's input value ${\bf m}$.
Design a protocol for achieving Multi-Valued Byzantine Broadcast, and prove that it is secure. Your protocol should be able to tolerate $f < n$ number of corruptions.
(Hint: generalize the Dolev-Strong protocol, and prove using the same framework we learned in class.)

\end{exercise}

\begin{answer} 
Write Solution Here.
\end{answer}



% =============================================================================
\begin{exercise}
{\bf (15 points)}    
Consider an execution of the Streamlet protocol where $n = 99$. Suppose strictly fewer than $n / 3$ nodes are corrupt. 
Let 
\begin{align*}
    \bot - e_1 - e_2 - \ldots
\end{align*}
denote a notarized chain starting with the genesis block $\bot$ followed by blocks 
with epoch numbers $e_1$, $e_2$, and so on.
For each of the following scenarios, please answer whether such a scenario can possibly occur during the execution. If so, please provide an explanation of how such a scenario can occur, and state which blocks are considered final by honest nodes. If not, please explain why. 
% Explain the chains
\begin{enumerate}[(a)]
    \item The two \emph{notarized} chains $\bot - 1 - 3 - 5 - 7$ and $\bot - 1 - 2 - 4 - 6$ will both be in honest view.

    \item An honest node $P$ observes a notarized chain of the form $\bot - 1 - 2 - 3 - 4$ and an honest node $Q$ observes a notarized chain of the form $\bot - 1 - 2 - 3 - 5$.

    \item An honest node $P$ observes a notarized chain of the form $\bot - 1 - 2 - 3 - 4 - 5$ and an honest node $Q$ observes a notarized chain of the form $\bot - 3 - 4 - 5$.
\end{enumerate}
\end{exercise}
\begin{answer}
Write Solution Here.
\end{answer}

% =============================================================================

\end{document}